% ---
% id: statistics/cor:tail-forms
% kind: cor
% title: "The two strict binomial tails"
% summary: >-
%   For every count k from 0 to n the strict tails of a binomial count N, the
%   probabilities that N lies below k and above k, are regularised incomplete
%   beta integrals, zero at k = 0 and at k = n respectively. These are the two
%   tails the endpoint rules compare with their budgets.
% prominence: standard
% status: implemented
% uses: [thm:binomial-tails]
% assumes: []
% refs: []
% provenance: newly-derived
% implements: certiqs_key.statistics.tails:upper_tail
% lean:
% ---

An endpoint rule compares the probability that the count falls strictly below
a lower endpoint, or strictly above an upper one, with a budget. Both strict
tails follow from Theorem~\ref{thm:binomial-tails} by a shift of the count.

\begin{corollary}\label{cor:tail-forms}
Let $n\ge1$, $p\in[0,1]$ and let $\Ncount{}$ be binomial with parameters
$(n,p)$, as in Theorem~\ref{thm:binomial-tails}. Then, for $0\le k\le n$,
\begin{equation}\label{eq:tails-used}
\begin{aligned}
  \Pr[\Ncount{}<k]&=\Pr[\Ncount{}\le k-1]
  =\begin{cases}0,&k=0,\\[2pt] \Ireg{1-p}{n-k+1}{k},&1\le k\le n,\end{cases}\\[4pt]
  \Pr[\Ncount{}>k]&=\Pr[\Ncount{}\ge k+1]
  =\begin{cases}\Ireg{p}{k+1}{n-k},&0\le k\le n-1,\\[2pt] 0,&k=n.\end{cases}
\end{aligned}
\end{equation}
\end{corollary}

\begin{proof}
Both are obtained from \eqref{eq:uppertail}--\eqref{eq:lowertail} by
substituting $k-1$ and $k+1$ respectively, $\Ncount{}$ being integer-valued.
For the first form: at $k=0$ the boundary value $\Pr[\Ncount{}\le k-1]=0$ for
$k-1<0$ of Theorem~\ref{thm:binomial-tails} applies; for $1\le k\le n$ one has
$0\le k-1\le n-1$, and \eqref{eq:lowertail} at $k-1$ gives
$\Ireg{1-p}{n-(k-1)}{(k-1)+1}=\Ireg{1-p}{n-k+1}{k}$. For the second form: for
$0\le k\le n-1$ one has $1\le k+1\le n$, and \eqref{eq:uppertail} at $k+1$
gives $\Ireg{p}{k+1}{n-(k+1)+1}=\Ireg{p}{k+1}{n-k}$; at $k=n$ the boundary
value $\Pr[\Ncount{}\ge k+1]=0$ for $k+1>n$ of Theorem~\ref{thm:binomial-tails} applies.
\end{proof}
